\scanpage{20}
Alors $G=T$, donc ici encore \deleted{\unsure{$G$}} $D_\sigma^*=D_\sigma$ est un $G$-torseur

\uline{Dém} Il est évident dans l'un et l'autre cas que \inserted{\text{$G_0={}$}} $k^*$ resp $T$ $(\simeq k^{+})$ commute à $\sigma$ et $\sigma'$, donc à $\Gamma$ (engendré par $\sigma$ et $\sigma'$) donc
\[
\deleted{k^*}G_0\subset G\qquad\deleted{\text{resp } T\subset G}.
\]
Or $\deleted{k^*}G_0$ est transitif sur $D_\sigma^*$ \deleted{\unsure{g}} (dans l'un et l'autre cas) et comme $G$ \deleted{est} opère librement sur $D_\sigma^*$ \inserted{\text{grâce à \unsure{prop 1 ou 2}}}, il en résulte ``formellement'' que $G$ opère transitivement -- donc $D_\sigma^*$ est un $G$ torseur -- et que $G_0=G$, cqfd
% archive sheet number: 19
